\documentclass[12pt,a4paper]{article}
\usepackage[margin=25mm, headheight=15pt, headsep=10pt]{geometry}
\usepackage{fontspec}
\usepackage[table]{xcolor} % Note: table option is required for rowcolors
\usepackage{titlesec}
\usepackage{tocloft}
\usepackage{fancyhdr}
\usepackage{booktabs}
\usepackage{tcolorbox}
\tcbuselibrary{skins, breakable}
\usepackage[colorlinks=true, linkcolor=emerald, urlcolor=emerald, bookmarks=true]{hyperref}

\definecolor{emerald}{RGB}{0,102,51}
\definecolor{darkred}{RGB}{139,0,0}
\definecolor{lightgray}{RGB}{245,245,245}

\usepackage[nil]{babel}
\babelprovide[import, main]{bengali}
\babelprovide[import]{arabic}
\babelfont{rm}[Renderer=HarfBuzz,Scale=1.0]{Noto Serif Bengali}
\babelfont{sf}[Renderer=HarfBuzz]{Noto Sans Bengali}
\babelfont{tt}[Renderer=HarfBuzz]{Noto Sans Bengali}
\babelfont[arabic]{rm}[Renderer=HarfBuzz,Scale=1.1]{Noto Naskh Arabic}

% Section styling
\titleformat{\section}{\Large\bfseries\color{emerald}}{\thesection.}{0.5em}{}
\titleformat{\subsection}{\large\bfseries\color{emerald!85!black}}{\thesubsection.}{0.5em}{}
\titleformat{\subsubsection}{\normalsize\bfseries\color{emerald!70!black}}{\thesubsubsection.}{0.5em}{}
\titlespacing*{\section}{0pt}{18pt}{8pt}

% Fancy Header/Footer setup
\pagestyle{fancy}
\fancyhf{} % clear all header and footer fields
\fancyhead[L]{\color{emerald}\small\textbf{\nouppercase{\leftmark}}}
\fancyfoot[L]{\scriptsize\color{black!50}{\lat{AI}}-সংকলিত, আলেম-যাচাইকৃত নয়}
\fancyfoot[C]{\color{emerald!70!black}\thepage}
\renewcommand{\headrulewidth}{0.4pt}
\renewcommand{\headrule}{\hbox to\headwidth{\color{emerald}\leaders\hrule height \headrulewidth\hfill}}
\renewcommand{\footrulewidth}{0pt}

% Custom boxes
% 1. Ayat/Hadith Box
\newtcolorbox{ayatbox}[1][]{
    enhanced, breakable, 
    colback=emerald!5, 
    colframe=emerald, 
    boxrule=0.5pt, 
    arc=3pt, 
    left=10pt, right=10pt, top=10pt, bottom=10pt,
    #1
}

% 2. Quote/Translation Box
\newtcolorbox{quotebox}[1][]{
    enhanced, breakable,
    frame hidden,
    colback=lightgray,
    borderline west={3pt}{0pt}{emerald},
    left=15pt, right=10pt, top=10pt, bottom=10pt,
    sharp corners,
    #1
}

% 3. AI-generated notice box
\newtcolorbox{warnbox}[1][]{
    enhanced, breakable,
    colback=red!4,
    colframe=darkred,
    boxrule=0.8pt,
    arc=3pt,
    left=10pt, right=10pt, top=10pt, bottom=10pt,
    #1
}

% Commands
\newcommand{\arab}[1]{\foreignlanguage{arabic}{#1}}
\newcommand{\kib}[1]{\textcolor{emerald}{\textbf{#1}}}

% Latin (English) fallback: Noto Serif Bengali has NO Latin glyphs
\newfontfamily\latfont[Renderer=HarfBuzz]{Noto Serif}
\newcommand{\lat}[1]{{\latfont #1}}

\setlength{\parskip}{5pt}
\setlength{\parindent}{0pt}

% Fix for missing bullet in Noto Serif Bengali
\renewcommand{\labelitemi}{$\bullet$}



\begin{document}
\begin{center}{\Large\bfseries\color{emerald} আল্লাহর রহমত ও তাওবা — \lat{02}-নব্বই-নয়-জন-হত্যাকারীর-তাওবা}\end{center}
\vspace{1em}
\section*{ঘটনা ২ — ৯৯ জনের হত্যাকারী: \lat{“}তাওবার পথে কে বাধা দেবে?\lat{”}}

\begin{warnbox}
 \textbf{গুরুত্বপূর্ণ নোটিশ (\lat{Important} \lat{Notice}):} এই কন্টেন্টটি \lat{AI} (কৃত্রিম বুদ্ধিমত্তা) দিয়ে প্রস্তুত ও সংকলিত। \lat{AI} কঠোর সূত্র-মান নীতি মেনে শুধু নির্ভরযোগ্য উৎস থেকে তথ্য সংগ্রহ করেছে — কেবল সহীহ/হাসান হাদিস ও সহীহ সনদের আসার রাখা হয়েছে, দুর্বল সনদ বাদ দেওয়া হয়েছে। তবে এটি কোনো আলেম বা মুহাদ্দিস কর্তৃক যাচাইকৃত নয়।
\end{warnbox}

\begin{quote}
\textbf{সম্পর্কিত ফাইল:} কুরআন ফাইল — ফুরকান ২৫:৭০; হাদিস ফাইল — হত্যাকারীর তাওবা কবুলের অধ্যায়।
\end{quote}

\medskip\hrule\medskip

\subsection*{১. সূত্র কার্ড (\lat{Source} \lat{Card})}

\begin{table}[h]\centering\small
\begin{tabular}{ll}
বিষয় & বিবরণ \\
\hline
\textbf{ঘটনা} & ৯৯ জন হত্যার পর তাওবার খোঁজ; এক পূজারীর \lat{“}না\lat{”}-তে ১০০তম হত্যা; তারপর এক আলেমের পথনির্দেশে তাওবা — অন্তরের ফেরা ও রহমতের ফেরেশতাদের বিজয় \\
\textbf{বর্ণনাকারী} & আবু সাঈদ খুদরি (রাঃ) \\
\textbf{গ্রন্থ ও নম্বর} & সহীহ মুসলিম ২৭৬৬; সহীহ বুখারি ৩৪৭০ \\
\textbf{অধ্যায়} & কিতাবুত তাওবাহ — \lat{“}হত্যাকারীর তাওবা কবুল, তার হত্যা যত বেশি হোক\lat{”} \\
\textbf{সনদের মান} & \textbf{সহীহ} (বুখারি-মুসলিম) \\
\textbf{স্থান ও সময়} & বনি ইসরাঈলের যুগের ঘটনা — নবীজি (সা.) সাহাবিদের কাছে বর্ণনা করেছেন \\
\hline
\end{tabular}
\end{table}

\medskip\hrule\medskip

\subsection*{২. সংক্ষিপ্ত পটভূমি (\lat{Background})}

এটি কোনো নির্দিষ্ট সাহাবির ঘটনা নয় — বনি ইসরাঈলের এক ব্যক্তির ঘটনা, যা নবীজি (সা.) সাহাবিদের শিখিয়েছেন। গুরুত্ব — ইসলামের ইতিহাসে সবচেয়ে বড় পাপ (ইচ্ছাকৃত হত্যা) করেও তাওবার দরজা খোলা থাকার দলিল। হাদিসটির অধ্যায়ের নামই বলে দেয়: \textbf{\lat{“}হত্যাকারীর তাওবা কবুল — তার হত্যা যত বেশি হোক।\lat{”}}

\medskip\hrule\medskip

\subsection*{৩. পূর্ণ ঘটনার বিশুদ্ধ বর্ণনা (\lat{The} \lat{Full} \lat{Narration})}

\begin{quote}
\textbf{\lat{“}তোমাদের আগে এক ব্যক্তি ছিল, যে নিরানব্বই জনকে হত্যা করেছিল। তারপর সে পৃথিবীর সবচেয়ে জ্ঞানী ব্যক্তির খোঁজ করল। তাকে এক পূজারীর (রাহিব) কাছে পাঠানো হলো। সে গিয়ে বলল: আমি নিরানব্বই জনকে হত্যা করেছি — আমার তাওবার কোনো সুযোগ আছে কি? পূজারী বলল: না। সে তাকেও হত্যা করল — এতে একশো হয়ে গেল। তারপর আবার পৃথিবীর সবচেয়ে জ্ঞানীর খোঁজ করল; তাকে এক আলেমের কাছে পাঠানো হলো। সে বলল: আমি একশো জনকে হত্যা করেছি — আমার তাওবার সুযোগ আছে কি? আলেম বললেন: হ্যাঁ! তাওবার পথে কে বাধা দেবে? তুমি অমুক জায়গায় যাও — সেখানে আল্লাহর ইবাদতকারী লোক আছে; তুমি তাদের সাথে আল্লাহর ইবাদত করো এবং নিজের জায়গায় ফিরে যেও না — কারণ তা মন্দ জায়গা। সে রওনা হলো — অর্ধেক পথ যেতেই তার মৃত্যু এলো। রহমতের ফেরেশতা ও শাস্তির ফেরেশতাদের মধ্যে বিবাদ হলো। রহমতের ফেরেশতারা বলল: সে তাওবাকারী হয়ে আল্লাহর দিকে অন্তর নিয়ে এসেছে। শাস্তির ফেরেশতারা বলল: সে তো কোনো ভালো কাজই করেনি। তখন এক ফেরেশতা মানুষের রূপে এসে মীমাংসা করলেন: দুই জায়গার মধ্যে দূরত্ব মাপো — যে জায়গার দিকে সে বেশি কাছে, তারই (সে)। মাপা হলো — সে ইবাদতের জায়গার বেশি কাছে ছিল; রহমতের ফেরেশতারা তাকে নিয়ে গেল।\lat{”}}
\end{quote}

\begin{quote}
কাতাদা বলেন — হাসান (বসরি) আমাদের বলেছেন: বর্ণনা করা হয়েছে — মৃত্যু যখন এলো, সে বুকে ভর করে হামাগুড়ি দিয়ে ইবাদতের জায়গার দিকে এগিয়ে গিয়েছিল।
\end{quote}

\medskip\hrule\medskip

\subsection*{৪. শব্দের অর্থ (\lat{Key} \lat{Words})}

\begin{table}[h]\centering\small
\begin{tabular}{ll}
শব্দ & অর্থ \\
\hline
\textbf{রাহিব (\arab{راهب})} & পূজারী/সন্ন্যাসী — এখানে ইবাদতে ডুবে থাকা ব্যক্তি \\
\textbf{আ'লামি আহলিল আরদ (\arab{أعلم} \arab{أهل} \arab{الأرض})} & পৃথিবীর সবচেয়ে জ্ঞানী — আলেম \\
\textbf{মান ইয়াহুলু বাইনাহু ওয়া বাইনাত তাওবাহ} & কে তার ও তাওবার মাঝে বাধা দাঁড়াবে? \\
\textbf{তাইবুন মুকবিলুন বিক্বলবিহি} & তাওবাকারী — অন্তর নিয়ে আল্লাহর দিকে ফেরা \\
\textbf{আরদু সাও' (\arab{أرض} \arab{سوء})} & মন্দ জায়গা — পাপের পরিবেশ \\
\hline
\end{tabular}
\end{table}

\medskip\hrule\medskip

\subsection*{৫. সংশ্লিষ্ট দলিল ও রেফারেন্স}

\begin{quote}
\textbf{\lat{“}তবে যে তাওবা করেছে, ঈমান এনেছে এবং সৎকাজ করেছে — তাদের পাপগুলো আল্লাহ নেকিতে পরিবর্তন করে দেবেন।\lat{”}} — সূরা আল-ফুরকান ২৫:৭০
\end{quote}

\begin{quote}
\textbf{\lat{“}তোমরা আল্লাহর রহমত থেকে নিরাশ হয়ো না। নিশ্চয়ই আল্লাহ সমস্ত গুনাহ ক্ষমা করে দেন।\lat{”}} — সূরা আয-যুমার ৩৯:৫৩
\end{quote}

\begin{quote}
\textbf{\lat{“}আল্লাহ রাতে তাঁর হাত বাড়ান… যতক্ষণ না পশ্চিম দিক থেকে সূর্য উদিত হয়।\lat{”}} — সহীহ মুসলিম ২৭৫৯
\end{quote}

\subsubsection*{মূল শিক্ষা (দলিলভিত্তিক)}

\begin{enumerate}
\item \textbf{হত্যার পরেও দরজা খোলা:} সবচেয়ে বড় পাপ — ১০০ জনের ইচ্ছাকৃত হত্যা — তারপরও আলেম বললেন \lat{“}হ্যাঁ, তাওবার পথে কে বাধা দেবে?\lat{”}; ইচ্ছাকৃত মুমিন হত্যার সতর্কবাণী (নিসা ৪:৯৩) সত্ত্বেও তাওবার দরজা আল্লাহ বন্ধ করেননি (ফুরকান ২৫:৭০-এর ব্যাপকতা — তাফসীর ইবনে কাসীর)।
\item \textbf{অন্তরের ফেরাই মূল:} সে কোনো নেক আমল করেনি — কিন্তু অন্তর আল্লাহর দিকে ফিরেছিল; রহমতের ফেরেশতারাই জিতেছে। আল্লাহ আমলের সংখ্যা নয়, অন্তরের দিক দেখেন।
\item \textbf{পূজারীর ভুল শিক্ষা:} \lat{“}তোমার তাওবা হবে না\lat{”} — এমন কথা কোনো আলেম বলতে পারে না; এতে মানুষকে হতাশ করাই শয়তানের কাজ। আলেমের ভাষা — \lat{“}তাওবার পথে কে বাধা দেবে?\lat{”}
\item \textbf{পরিবেশ বদল তাওবার অঙ্গ:} \lat{“}মন্দ জায়গা ছেড়ে ভালো জায়গায় যাও\lat{”} — তাওবাকে টিকিয়ে রাখতে পরিবেশ বদল জরুরি (ব্যবহারিক গাইড ৭ নং দেখুন)।
\end{enumerate}
\medskip\hrule\medskip

\subsection*{৬. রেফারেন্স}

\begin{itemize}
\item সহীহ মুসলিম ২৭৬৬ (কিতাবুত তাওবাহ — বাবু কুবুলি তাওবাতিল কাতিল)
\item সহীহ বুখারি ৩৪৭০ (কিতাবুল আম্বিয়া)
\item কুরআন: ফুরকান ২৫:৭০, যুমার ৩৯:৫৩, নিসা ৪:৯৩
\item ব্যাখ্যা: শারহু নববী (সহীহ মুসলিম), ফাতহুল বারি — হত্যাকারীর তাওবার ব্যাখ্যা
\end{itemize}

\end{document}
